\rtmtight
\begin{tblr}{width=\textwidth,colspec={|Q[c,m,2.1cm]|X[l,m]|},hlines,vlines,hspan=minimal,rowsep=1.5pt,colsep=3pt,row{1}={bg=headgray,font=\bfseries}}
\SetCell{c,m}\hfil\logo\hfil & \textbf{POLITEKNIK NEGERI MALANG}\par\textbf{JURUSAN TEKNOLOGI INFORMASI}\par\textbf{PROGRAM STUDI: D4 TEKNIK INFORMATIKA} \\
\SetCell[c=2]{l} \textbf{RENCANA TUGAS MAHASISWA (RTM) DAN RUBRIK PENILAIAN} & \\
Nama Mata Kuliah & Cloud Computing \\
Kode Mata Kuliah & RTI256107 \quad \textbf{sks / Jam perminggu:} 2 SKS/4 jam \quad \textbf{Semester:} 6 \\
Dosen Pengampu & Tim Pengajar Program Studi D4 TI \\
\end{tblr}
\assessmentblock{Tugas Arsitektur Cloud}{Case Method dan praktik cloud}{SCPMK0509-05101, SCPMK0905-05102, SCPMK1003-05103}{Mahasiswa merancang dan mendokumentasikan arsitektur cloud yang scalable dengan mempertimbangkan model layanan, keamanan, dan biaya sesuai studi kasus yang diberikan.}{Menganalisis kebutuhan, merancang arsitektur cloud menggunakan IaaS/PaaS/SaaS, mendokumentasikan keputusan teknis, dan mempresentasikan solusi.}{Diagram arsitektur cloud, laporan teknis, estimasi biaya, dan/atau bahan presentasi.}{Kelengkapan dan ketepatan rancangan arsitektur cloud & 12 & Arsitektur mencakup semua komponen (compute, storage, network, security), pemilihan model layanan tepat sesuai kebutuhan, dan dilengkapi justifikasi teknis yang kuat. & Arsitektur mencakup komponen utama tetapi beberapa justifikasi pemilihan layanan atau komponen belum lengkap. & Arsitektur tidak lengkap, pemilihan layanan tidak sesuai kebutuhan, atau tidak terdokumentasi. \\ \\
Implementasi containerisasi dan model virtualisasi & 10 & Konsep virtualisasi, Docker, dan Kubernetes diterapkan dengan benar dalam rancangan; Dockerfile dan manifest Kubernetes valid dan sesuai best practice. & Containerisasi diimplementasikan tetapi konfigurasi Docker atau Kubernetes masih memiliki kekurangan teknis. & Containerisasi tidak diterapkan atau implementasi menghasilkan konfigurasi yang tidak fungsional. \\ \\
Analisis solusi cloud berdasarkan keandalan dan biaya & 8 & Solusi dievaluasi berdasarkan SLA, keandalan, skalabilitas, dan estimasi biaya menggunakan pricing calculator platform cloud. & Evaluasi solusi ada tetapi analisis biaya atau parameter keandalan belum komprehensif. & Tidak ada analisis solusi yang memadai atau estimasi biaya tidak realistis. \\}

\assessmentblock{Kuis Konsep dan Evaluasi Cloud}{Kuis}{SCPMK0905-05102, SCPMK1003-05103}{Kuis untuk mengukur pemahaman konsep virtualisasi, model layanan cloud, dan kemampuan mengevaluasi solusi cloud.}{Menjawab soal pilihan ganda dan/atau uraian terkait konsep cloud computing, virtualisasi, containerisasi, keamanan, dan jaringan cloud.}{Jawaban tertulis kuis.}{Pemahaman konsep virtualisasi dan model layanan cloud & 5 & Konsep hypervisor, VM, container, IaaS/PaaS/SaaS dijelaskan dengan benar dan perbedaan antar model layanan dipahami secara mendalam. & Sebagian besar konsep dipahami tetapi penjelasan perbedaan model layanan atau contoh penerapan masih kurang. & Konsep tidak dipahami atau terdapat kekeliruan mendasar dalam penjelasan model layanan. \\ \\
Kemampuan evaluasi dan pemilihan solusi cloud & 5 & Solusi cloud dievaluasi dengan tepat berdasarkan kriteria keandalan, keamanan, skalabilitas, dan biaya; rekomendasi didukung argumen teknis. & Evaluasi sebagian besar tepat tetapi beberapa kriteria evaluasi atau rekomendasi belum kuat. & Evaluasi tidak tepat atau tidak menggunakan kriteria yang relevan. \\}

\assessmentblock{UTS Virtualisasi dan Arsitektur Cloud}{UTS berbasis pemrograman dan desain}{SCPMK0905-05102, SCPMK0509-05101}{Evaluasi tengah semester untuk mengukur kemampuan menjelaskan virtualisasi dan merancang arsitektur cloud dasar.}{Menjawab soal uraian dan/atau praktik terkait virtualisasi, containerisasi Docker, dan rancangan arsitektur IaaS/PaaS.}{Jawaban tertulis, konfigurasi Docker/cloud, dan dokumen rancangan arsitektur.}{Penguasaan konsep virtualisasi dan containerisasi & 7 & Konsep hypervisor, VM, dan container dijelaskan dengan tepat; Dockerfile valid dan konfigurasi container sesuai kebutuhan. & Sebagian besar konsep dijelaskan benar tetapi konfigurasi container atau detail teknis masih ada kekurangan. & Konsep virtualisasi tidak dipahami atau konfigurasi container tidak berfungsi. \\ \\
Rancangan arsitektur IaaS dan PaaS & 6 & Arsitektur IaaS/PaaS dirancang dengan komponen yang tepat, relasi antar layanan jelas, dan sesuai best practice platform cloud. & Rancangan arsitektur ada tetapi beberapa komponen atau relasi layanan belum optimal. & Rancangan arsitektur tidak mencerminkan pemahaman IaaS/PaaS atau tidak dapat diimplementasikan. \\}

\assessmentblock{PBL Deployment Aplikasi Cloud}{Project Based Learning}{SCPMK0509-05101, SCPMK1003-05103}{Mahasiswa mendeployment aplikasi nyata ke platform cloud menggunakan container, mengkonfigurasi keamanan dan jaringan, serta mengoptimalkan biaya operasional.}{Merancang arsitektur deployment, mengimplementasikan containerisasi, mengkonfigurasi VPC dan IAM, menguji fungsionalitas, dan mendokumentasikan.}{Aplikasi yang berjalan di cloud, konfigurasi infrastruktur (IaC), laporan teknis, dan bukti pengujian.}{Keberhasilan deployment dan konfigurasi infrastruktur & 9 & Aplikasi berhasil di-deploy ke platform cloud dengan containerisasi yang benar, networking (VPC/Load Balancer) terkonfigurasi, dan endpoint dapat diakses. & Deployment berhasil tetapi konfigurasi networking atau containerisasi masih memiliki kekurangan teknis. & Deployment gagal atau aplikasi tidak dapat diakses dari luar environment cloud. \\ \\
Implementasi keamanan dan IAM & 7 & IAM dikonfigurasi dengan prinsip least privilege, enkripsi data-at-rest dan in-transit diterapkan, dan kepatuhan terhadap best practice keamanan terdokumentasi. & Konfigurasi IAM dan enkripsi ada tetapi prinsip least privilege atau dokumentasi keamanan belum lengkap. & IAM tidak dikonfigurasi atau keamanan dasar tidak diterapkan. \\ \\
Optimasi biaya dan dokumentasi & 6 & Strategi optimasi biaya diterapkan (right-sizing, auto-scaling), estimasi biaya bulanan realistis, dan dokumentasi infrastruktur lengkap. & Optimasi biaya dipertimbangkan tetapi strategi atau dokumentasi belum komprehensif. & Tidak ada pertimbangan optimasi biaya atau dokumentasi infrastruktur tidak memadai. \\}

\assessmentblock{UAS Evaluasi Arsitektur Cloud dan Presentasi}{UAS praktik dan presentasi}{SCPMK0509-05101, SCPMK1003-05103}{Mahasiswa mengevaluasi dan mempresentasikan rancangan arsitektur cloud komprehensif yang mencakup skalabilitas, keamanan, jaringan, dan strategi multi-cloud.}{Merancang arsitektur cloud end-to-end, melakukan analisis komparatif platform cloud, mendokumentasikan trade-off, dan mempresentasikan defense teknis.}{Dokumen arsitektur cloud, analisis perbandingan platform, laporan evaluasi, dan slide presentasi.}{Kelengkapan arsitektur cloud end-to-end & 10 & Arsitektur mencakup semua lapisan (compute, storage, network, security, monitoring), scalability dan disaster recovery dirancang, serta semua komponen terdokumentasi lengkap. & Arsitektur mencakup sebagian besar lapisan tetapi beberapa aspek seperti DR atau monitoring belum dirancang secara detail. & Arsitektur tidak lengkap atau tidak mempertimbangkan aspek keandalan dan skalabilitas. \\ \\
Analisis komparatif dan pemilihan platform cloud & 8 & Platform cloud (AWS/GCP/Azure) dibandingkan secara objektif berdasarkan fitur, SLA, biaya, dan kebutuhan organisasi; rekomendasi didukung data dan argumen teknis. & Perbandingan platform dilakukan tetapi analisis biaya, SLA, atau justifikasi rekomendasi belum kuat. & Tidak ada analisis komparatif yang memadai atau rekomendasi tidak berdasarkan kriteria teknis. \\ \\
Kualitas defense dan argumentasi teknis & 7 & Mahasiswa mampu mempertahankan setiap keputusan arsitektur dengan evidence teknis, menjelaskan trade-off, dan merespons pertanyaan penguji secara meyakinkan. & Defense cukup baik tetapi beberapa keputusan arsitektur tidak dapat dipertahankan dengan evidence yang kuat. & Tidak mampu menjelaskan keputusan arsitektur atau tidak dapat merespons pertanyaan teknis. \\}

\begin{tblr}{width=\textwidth,colspec={|Q[l,m,3.2cm]|X[l,m]|},hlines,vlines,hspan=minimal,row{1-3}={bg=headgray},rowsep=1.5pt,colsep=3pt}
Jadwal Pelaksanaan: & Minggu 1--17 sesuai RPS. Setiap evaluasi memiliki RTM dan rubrik multi-kriteria. \\
Lain-Lain Yang Diperlukan: & LMS, akun cloud (AWS/GCP/Azure), Docker, Kubernetes, perangkat lunak sesuai mata kuliah, dan perangkat presentasi. \\
Pustaka: & Mengacu pustaka utama dan pendukung pada bagian RPS. \\
\end{tblr}
